\section*{第 \ref{ch:g1:the-five-senses} 章——五种感觉}

\begin{solution}{exo:g1:the-five-senses:1}
视觉——眼睛；听觉——耳朵；嗅觉——鼻子；
味觉——舌头；触觉——\omterm{def:g1:the-five-senses:sense}{皮肤}。
\end{solution}

\begin{solution}{exo:g1:the-five-senses:2}
(a) 眼睛。(b) \omterm{def:g1:the-five-senses:sense}{皮肤}。(c) 鼻子。(d) 耳朵。
\end{solution}

\begin{solution}{exo:g1:the-five-senses:3}
\omterm{def:g1:the-five-senses:sense}{皮肤}，是触觉器官。它覆盖全身，
从头到脚。
\end{solution}

\begin{solution}{exo:g1:the-five-senses:4}
手上的\omterm{def:g1:the-five-senses:sense}{皮肤}传来信息（冷、硬、凹凸不平）；
大脑理解了这一信息，并说出了物体的名字。
\end{solution}

\begin{solution}{exo:g1:the-five-senses:5}
味觉和嗅觉。鼻子承担了风味的一半工作，
所以鼻子被堵住时，食物就显得没有味道。
\end{solution}

\begin{solution}{exo:g1:the-five-senses:6}
例如：千万不要直视太阳——保护眼睛；
不让声音过大——保护耳朵；吹凉热食——保护舌头；
不把任何东西塞进耳朵或鼻子——保护耳朵和鼻子；
读书时光线要好——保护眼睛。
\end{solution}

\begin{solution}{exo:g1:the-five-senses:7}
因为声音太大会伤害耳朵，而受伤的耳朵
可能会一直受伤。机器罩能让机器噪声足够小，
使人一生都有好的听力。
\end{solution}

\begin{solution}{exo:g1:the-five-senses:8}
开放性答案——多数人能说出三种中的两种或三种。
鼻子能很好地识别熟悉的食物，
即使不用眼睛帮忙。
\end{solution}

\begin{solution}{exo:g1:the-five-senses:9}
眼睛只是\emph{看见}——它收集光线并传递信息。
大脑负责理解：它接收信息，并说出“这是奶奶”、
“这是我的学校”。感觉负责收集，
大脑负责理解。
\end{solution}

\begin{solution}{exo:g1:the-five-senses:10}
阅读时用指尖——触觉代替视觉，
摸读纸上凸起的盲文圆点。过马路时用耳朵——听汽车来来往往的声音——
常常还会拿着盲杖，
用杖尖的触觉探查前方的地面。
\end{solution}
